\appendix
\chapter{实验环境}
% TODO: 草稿附录“实验环境/硬件配置/软件工具”仍为模板占位，需要补充真实内容。

\section{研究记录}
\subsection{2025.10.17}
完成内容：数据分析项目框架，使用 python 语言搭建，包含主运行函数，各个数据处理函数（dataio），将 tdms 以及 tdms\_index 文件转换成 csv 文件和 json 元数据文件；将时域 csv 文件进行 FFT 处理，生成频域 csv 文件；将每个实验组的频域 csv 文件进行平均，得到频域平均 csv 文件（目前的数据不同文件之间的时间点是间断的，无法进行合并分析、增量分析绘制 RGB 图）。

接下来内容：1、使用comet等，查找与本研究方法相似的相关数据处理论文，参考其数据处理方法，寻找优化的数据处理方法。2、对于频域平均后的csv文件，分析所有可能的分析方法。3、对比一个实验组中的20个文件，分析其频域、时域上是否具有足够的稳定性，为频域平均提供理论支撑。
